%% start of file `template.tex'.
%% Copyright 2006-2013 Xavier Danaux (xdanaux@gmail.com).
%
% This work may be distributed and/or modified under the
% conditions of the LaTeX Project Public License version 1.3c,
% available at http://www.latex-project.org/lppl/.


\documentclass[11pt,a4paper,roman]{moderncv}        % possible options include font size ('10pt', '11pt' and '12pt'), paper size ('a4paper', 'letterpaper', 'a5paper', 'legalpaper', 'executivepaper' and 'landscape') and font family ('sans' and 'roman')

% modern themes
\moderncvstyle{banking}                            % style options are 'casual' (default), 'classic', 'oldstyle' and 'banking'
\moderncvcolor{blue}                                % color options 'blue' (default), 'orange', 'green', 'red', 'purple', 'grey' and 'black'
%\renewcommand{\familydefault}{\sfdefault}         % to set the default font; use '\sfdefault' for the default sans serif font, '\rmdefault' for the default roman one, or any tex font name
\nopagenumbers{}                                  % uncomment to suppress automatic page numbering for CVs longer than one page

% character encoding
\usepackage[utf8]{inputenc}
\usepackage{fontawesome}
\usepackage{academicons}
% \usepackage{fontawesome5} % for \faGoogleScholar
\usepackage{tabularx}
\usepackage{ragged2e}

% if you are not using xelatex ou lualatex, replace by the encoding you are using
%\usepackage{CJKutf8}                              % if you need to use CJK to typeset your resume in Chinese, Japanese or Korean

% adjust the page margins
\usepackage[scale=0.8]{geometry}
\usepackage{multicol}
%\setlength{\hintscolumnwidth}{3cm}                % if you want to change the width of the column with the dates
%\setlength{\makecvtitlenamewidth}{10cm}           % for the 'classic' style, if you want to force the width allocated to your name and avoid line breaks. be careful though, the length is normally calculated to avoid any overlap with your personal info; use this at your own typographical risks...

\usepackage{import}
\definecolor{honor}{RGB}{197,48,48}   % Oral / Spotlight / best-paper marks (matches website accent)

% Publication list: bold venue tag in a hanging-indent left column.
\newenvironment{publist}{%
  \begin{list}{}{%
    \setlength{\leftmargin}{34mm}%
    \setlength{\labelwidth}{32mm}%
    \setlength{\labelsep}{2mm}%
    \setlength{\itemsep}{1.4mm}%
    \setlength{\parsep}{0pt}%
    \setlength{\topsep}{1.5mm}%
    \renewcommand{\makelabel}[1]{\bfseries ##1\hfill}%
  }}{\end{list}}

% personal data
\name{Shenglai}{Zeng}
% \title{Curriculum Vitae}                               % optional, remove / comment the line if not wanted
% optional, remove / comment the line if not wanted; the "postcode city" and and "country" arguments can be omitted or provided empty
 %\phone[mobile]{(+88)0 1718268856}                   % optional, remove / comment the line if not wanted
% \phone[fixed]{01234 123456}                    % optional, remove / comment the line if not wanted
%\phone[fax]{+3~(456)~789~012}                      % optional, remove / comment the line if not wanted
% \email{xpan1@swarthmore.edu}                               % optional, remove / comment the line if not wanted
% \homepage{shawnpan.me}                         % optional, remove / comment the line if not wanted
% \extrainfo{}                 % optional, remove / comment the line if not wanted
%\photo[64pt][0.4pt]{picture}                       % optional, remove / comment the line if not wanted; '64pt' is the height the picture must be resized to, 0.4pt is the thickness of the frame around it (put it to 0pt for no frame) and 'picture' is the name of the picture file
%\quote{Some quote}                                 % optional, remove / comment the line if not wanted

% to show numerical labels in the bibliography (default is to show no labels); only useful if you make citations in your resume
%\makeatletter
%\renewcommand*{\bibliographyitemlabel}{\@biblabel{\arabic{enumiv}}}
%\makeatother
%\renewcommand*{\bibliographyitemlabel}{[\arabic{enumiv}]}% CONSIDER REPLACING THE ABOVE BY THIS

% bibliography with mutiple entries
%\usepackage{multibib}
%\newcites{book,misc}{{Books},{Others}}
  
\newcommand*{\customcventry}[7][.25em]{
  \begin{tabular}{@{}l} 
    {\bfseries #4}
  \end{tabular}
  \hfill% move it to the right
  \begin{tabular}{l@{}}
     {\bfseries #5}
  \end{tabular} \\
  \begin{tabular}{@{}l} 
    {\itshape #3}
  \end{tabular}
  \hfill% move it to the right
  \begin{tabular}{l@{}}
     {\itshape #2}
  \end{tabular}
  \ifx&#7&%
  \else{\\%
    \begin{minipage}{\maincolumnwidth}%
      \small#7%
    \end{minipage}}\fi%
  \par\addvspace{#1}}

\newcommand*{\customcvproject}[4][.25em]{
%   \vfill\noindent
  \begin{tabular}{@{}l} 
    {\bfseries #2}
  \end{tabular}
  \hfill% move it to the right
  \begin{tabular}{l@{}}
     {\itshape #3}
  \end{tabular}
  \ifx&#4&%
  \else{\\%
    \begin{minipage}{\maincolumnwidth}%
      \small#4%
    \end{minipage}}\fi%
  \par\addvspace{#1}}

\setlength{\tabcolsep}{12pt}

%----------------------------------------------------------------------------------
%            content
%----------------------------------------------------------------------------------
\begin{document}
%\begin{CJK*}{UTF8}{gbsn}                          % to typeset your resume in Chinese using CJK
%-----       resume       ---------------------------------------------------------
\makecvtitle
\vspace*{-23mm}

% \begin{center}

% \begin{tabular}{ c c c c }
%  \\
%  \\
%  \faGlobe\enspace https://slz-ai.github.io/&
%  \faEnvelopeO\enspace shenglaizeng@gmail.com & \faMobile\enspace (+1) 517-974-7616 \\  
% \end{tabular}
% \end{center}

\begin{center}
\begin{tabular}{ c c c c }
\\
\\
\faGlobe\enspace \href{https://slz-ai.github.io/}{Personal Website}& % 使用 \href 包裹主页链接
\faGraduationCap \enspace \href{https://scholar.google.com/citations?user=2EFS5WEAAAAJ}{Google Scholar} & % **新增的 Google Scholar 链接**
\faEnvelopeO\enspace zengshe1@msu.edu & \faMobile\enspace (+1) 517-974-7616 \\
\end{tabular}
\end{center}

\section{RESEARCH INTEREST}
I am a fourth-year Ph.D. candidate at Michigan State University, advised by Professor Jiliang Tang, and I am on the job market for positions starting in 2027. My research builds \textbf{personalized, proactive, and safe agents} with \textbf{multimodal capabilities} that interact with people. I work on the two mechanisms that determine what such an agent can actually do: the \textbf{agent harness} (memory, retrieval, self-improvement, and evaluation environments) and \textbf{post-training}. I have published 35+ papers at NeurIPS, ICML, ICLR, ACL, EMNLP, NAACL, KDD, and CVPR, received the Best Paper Award of IEEE Transactions on Cloud Computing (2023) and an ACL-2026 SAC Highlight (top 1\%, Best Paper Candidate), and have industry research experience at Meta AI, Amazon, and Baidu.

\section{EDUCATION}

{\customcventry{Sept 2023-Present}{
DSE Lab/Ph.D. Candidate in Computer Science and Engineering   }{Michigan State University}{East Lansing, U.S}{}{}}
Expected Graduation: May 2027\\
Advisor: Jiliang Tang\\
Lab: Data Science and Engineering Lab\\
Research Focus: Agentic AI, Retrieval \& Memory, Personalized LLMs, Agent Safety \& Privacy, Multimodal

{\customcventry{Sept 2019 - Jun 2023}{
Yingcai Honor School/B.Sc in Computer Science and Engineering   }{University of Electronic Science and Technology of China }{Chengdu, China}{}{}}

CGPA: 3.98/4.00  \\
Weighted Average: 93.97/100(1st among 100 students)\\
Honors: The Most Outstanding Students Award of UESTC



%{\customcventry{August 2011}{Bachelor of Technology in Applied Electronics and Instrumentation Engineering   }{West Bengal University of Technology}{Kolkata, West India}{}{}}
\section{HONORS AND AWARDS}
\begin{minipage}{\maincolumnwidth}%
	\small{
    	\begin{itemize}
          \item \textbf{SAC Highlight (Top 1\%) \& Best Paper Candidate}, ACL-2026 \hfill 2026
          \item \textbf{Best Paper Award}, IEEE Transactions on Cloud Computing (JCR Q1) \hfill 2023
          \item \textbf{The Most Outstanding Students Award}, UESTC (highest honor; 10 students university-wide)
          \item \textbf{National Scholarship} (highest undergraduate honor in China) \hfill 2020
          \item WAC Scholarship (10 undergraduates university-wide per year), UESTC \hfill 2021
          \item 1st Outstanding Academic Scholarship (Top 5\%), UESTC \hfill 2020, 2021, 2022
		\end{itemize}}%
\end{minipage}%

\section{INDUSTRY EXPERIENCE}
\textit{All papers in this section are \textbf{first-authored} (including co-first).}

{\customcventry{Summer 2026}{Research Scientist Intern/Next Generation Proactive Agents}{Meta AI}{Seattle, USA}{}{}
-Mentored by Xiao Yang, \href{https://lunadong.com/}{Luna Dong}, and \href{https://scottyih.org/}{Scott Yih}\\
-Closely worked with a cross-organizational team spanning Superintelligence Labs, Reality Labs, FAIR, and MRS\\
I built \textbf{memory-driven proactive agents}: assistants that can proactively use their context and decide \emph{whether and when} to speak, not just what to answer. I introduced \textbf{MPA}, a benchmark pairing egocentric video with year-long, distractor-rich personal memories (\textbf{$\sim$618K tokens} each) and scoring content and timing jointly, and found \textbf{existing models severely limited} in this regime. My plug-and-play harness \textbf{GPM} surfaces memory from a graph-organized store and gates it down to what is worth interrupting for, lifting strict accuracy from \textbf{less than 10\%} to \textbf{37.0\%} (Qwen3.8-27B) and to \textbf{38\%} on GPT-5.3-Chat with no training at all, reaching \textbf{42.8\%} with \textbf{RL post-training}.
\begin{itemize}
    \item Memory-Driven Proactive Agents \textbf{[Submitted to ICLR'27]}.
\end{itemize}
}

{\customcventry{Summer 2024, 2025}{Applied Scientist Intern/Retrieval-augmented Generation and Personalization}{Amazon, Search \& Shopping Foundation Model Team}{CA, USA}{}{}
-Mentor: Tianqi Zheng, Dante Everaert \\
-Manager: Hanqing Lu, Chuan Tian, Monica Xiao Cheng
\\[0.4em]
My work centered on two sides of making retrieval-augmented and personalized LLMs deployable. \textbf{On the model side} (how a system uses the context it retrieves): attention-guided compression cuts personalization context \textbf{50$\times$} (10K$\rightarrow$200 tokens) at near-full-context quality \href{https://arxiv.org/abs/2602.07778}{[ACL'26]}; gated representation fine-tuning buys \textbf{+26--35 points} of context robustness while tuning only \textbf{0.0004\% of parameters on fewer than 200 examples} \href{https://arxiv.org/abs/2502.14100}{[ACL'25]}; and representation-based knowledge checking filters misleading evidence at \textbf{95\% accuracy} \href{https://arxiv.org/abs/2411.14572}{[NAACL'25]}. \textbf{On the data side} (what the retrieval corpus itself exposes): \textbf{self-improving synthetic data} replaces private records at comparable task quality \href{https://arxiv.org/abs/2406.14773}{[EMNLP'25]}, and I gave the first account of how \textbf{multimodal RAG leaks} the images and text it retrieves \href{https://arxiv.org/abs/2505.13957}{[EMNLP'25]}. Invited back for a second summer.
\begin{itemize}
    \item Efficient Personalized LLMs (PLLMs) \href{https://arxiv.org/abs/2602.07778}{[ACL'26, Oral, Best Paper Candidate, SAC Highlight]}.
    \item Knowledge Checking in RAG \href{https://arxiv.org/abs/2411.14572}{[NAACL'25, Oral]}.
    \item Representation Finetuning for Robust RAG \href{https://arxiv.org/abs/2502.14100}{[ACL'25]}.
    \item Self-Improving Synthetic Data for RAG \href{https://arxiv.org/abs/2406.14773}{[EMNLP'25]}.
    \item Multi-modal RAG \href{https://arxiv.org/abs/2505.13957}{[EMNLP'25]}.
\end{itemize}
}

{\customcventry{May 2023 - May 2024}{Research Intern/RAG Privacy and LLM Memorization}{Baidu Inc., Search \& Agent Business Team}{Beijing, China}{}{}
-Mentor: Dr. Yiding Liu and Dr. Dawei Yin
\\[0.4em]
I established the privacy foundations of retrieval-augmented generation: the \textbf{first work to reveal that retrieval-augmented systems carry severe privacy risk}, leaking their own retrieval database \href{https://arxiv.org/abs/2402.16893}{[ACL'24]}. The threat model it introduced has since been carried into the privacy analysis of future systems (multi-modal RAG, agent memory, and multi-agent systems), and the paper now has \textbf{350+ citations}. I also produced the \textbf{first task-level map of memorization in LLM post-training} \href{https://arxiv.org/abs/2310.06714}{[ACL'24]}, showing how memorization \textbf{scales with model size} and that the scaling itself is task-dependent (22.3\% for summarization vs.\ 0.15\% for reading comprehension).
\begin{itemize}
    \item RAG Privacy and Retrieval-Database Leakage \href{https://arxiv.org/abs/2402.16893}{[ACL'24]}.
    \item Memorization in LLM post-training \href{https://arxiv.org/abs/2310.06714}{[ACL'24]}.
\end{itemize}
}



\section{PUBLICATIONS}
35+ papers at top-tier ML/NLP venues (NeurIPS, ICML, ICLR, ACL, EMNLP, NAACL, KDD, CVPR), including \textbf{12 first-authored} publications. \textbf{1340+ citations} (Google Scholar, Sept.\ 2026). Recognitions include a \textbf{Best Paper Award} (IEEE TCC) and \textbf{Best Paper Candidate} (ACL), \textbf{Spotlight} (ICLR), \textbf{Oral} presentations (ACL, NAACL), and adoption in frontier-model technical reports (Muse Spark, Qwen, Kimi). \textit{* denotes equal contribution.}

\vspace{2mm}
{\large\bfseries First-Authored Publications}
\vspace{-1mm}
\begin{publist}
\item[{[NeurIPS'26]}] \textbf{S. Zeng}, Q. Wang, K. Guo, X. Dai, X. Long, H. Liu, ``Magnifying What Matters: Attention-Guided Adaptive Rendering for Visual Text Comprehension'', \textit{Conference on Neural Information Processing Systems}, 2026.
\item[{[ICML'26]}] \textbf{S. Zeng*}, J. Zhang*, K. Guo, X. Dai, H. Liu, J. Tang, Y. Chang, ``Fix Before Search: Benchmarking Agentic Query Visual Pre-processing in Multimodal Retrieval-Augmented Generation'', \textit{International Conference on Machine Learning}, 2026.
\item[{[ACL'26, Oral]}] \textbf{S. Zeng}, T. Zheng, C. Tian, D. Everaert, Y.-S. Wang, Y. Huang, M. J. Morais, R. Patki, J. Tian, X. Dai, K. Guo, M. X. Cheng, H. Liu, ``Attn-GS: Attention-Guided Context Compression for Efficient Personalized LLMs'', \textit{Annual Meeting of the Association for Computational Linguistics, \textbf{\color{honor}Oral}, \textbf{\color{honor}SAC Highlight (Top 1\%)}, \textbf{\color{honor}Best Paper Candidate}}, 2026.
\item[{[EMNLP'25]}] \textbf{S. Zeng}, J. Zhang, P. He, J. Ren, T. Zheng, H. Lu, H. Xu, H. Liu, Y. Xing, J. Tang, ``Mitigating the Privacy Issues in Retrieval-Augmented Generation (RAG) via Pure Synthetic Data'', \textit{Conference on Empirical Methods in Natural Language Processing}, 2025.
\item[{[EMNLP'25]}] \textbf{S. Zeng*}, J. Zhang*, J. Ren, T. Zheng, H. Liu, X. Tang, Y. Chang, ``Beyond Text: Unveiling Privacy Vulnerabilities in Multi-modal Retrieval-Augmented Generation'', \textit{Conference on Empirical Methods in Natural Language Processing}, 2025.
\item[{[ACL'25]}] \textbf{S. Zeng}, P. He, K. Guo, T. Zheng, H. Lu, Y. Xing, H. Liu, ``Towards Context-Robust LLMs: A Gated Representation Fine-tuning Approach'', \textit{Annual Meeting of the Association for Computational Linguistics}, 2025.
\item[{[NAACL'25, Oral]}] \textbf{S. Zeng}, J. Zhang, B. Li, Y. Lin, T. Zheng, D. Everaert, H. Lu, H. Liu, Y. Xing, M. X. Cheng, J. Tang, ``Towards Knowledge Checking in Retrieval-Augmented Generation: A Representation Perspective'', \textit{Conference of the North American Chapter of the Association for Computational Linguistics, \textbf{\color{honor}Oral}}, 2025.
\item[{[ACL'24]}] \textbf{S. Zeng*}, Y. Li*, J. Ren, Y. Liu, H. Xu, P. He, Y. Xing, S. Wang, J. Tang, D. Yin, ``Exploring Memorization in Fine-tuned Language Models'', \textit{Annual Meeting of the Association for Computational Linguistics}, 2024.
\item[{[ACL'24]}] \textbf{S. Zeng*}, J. Zhang*, P. He, Y. Xing, Y. Liu, H. Xu, J. Ren, S. Wang, D. Yin, Y. Chang, J. Tang, ``The Good and The Bad: Exploring Privacy Issues in Retrieval-Augmented Generation (RAG)'', \textit{Findings of the Association for Computational Linguistics: ACL}, 2024.
\item[{[IEEE TCC, Best Paper]}] \textbf{S. Zeng}, Z. Li, H. Yu, Z. Zhang, L. Luo, B. Li, D. Niyato, ``HFedMS: Heterogeneous Federated Learning with Memorable Data Semantics in Industrial Metaverse'', \textit{IEEE Transactions on Cloud Computing, \textbf{\color{honor}Best Paper Award}}, 2023.
\item[{[DASFAA'22, Oral]}] \textbf{S. Zeng}, Z. Li, H. Yu, Y. He, Z. Xu, D. Niyato, H. Yu, ``Heterogeneous Federated Learning via Grouped Sequential-to-Parallel Training'', \textit{International Conference on Database Systems for Advanced Applications, \textbf{\color{honor}Oral}}, 2022.
\item[{[SDM'23]}] \textbf{S. Zeng*}, J. Wang*, Z. Long, Y. Wang, H. Xiao, F. Ma, ``Knowledge-Enhanced Semi-Supervised Federated Learning for Aggregating Heterogeneous Lightweight Clients in IoT'', \textit{SIAM International Conference on Data Mining}, 2023.
\end{publist}

\vspace{2mm}
{\large\bfseries Co-Authored Publications}
\vspace{-1mm}
\begin{publist}
\item[{[NeurIPS'26]}] Y. Li, Y. Feng, Z. Xu, Z. Ma, K. Zheng, F. Jiang, X. Sun, R. Shao, Z. Chen, Y. Huang, X. Han, B. Lee, K. Xu, \textbf{S. Zeng}, H. Hua, X. Zhang, B. Alomair, R. Krishna, L. Zettlemoyer, P. W. Koh, B. Ramasubramanian, L. Niu, X. Yue, R. Poovendran, ``JobBench: Aligning Agent Work With Human Will'', \textit{Conference on Neural Information Processing Systems, Datasets \& Benchmarks Track}, 2026. \textbf{(adopted by Meta Muse Spark 1.1, Moonshot Kimi K3, and Alibaba Qwen 3.8 Max)}
\item[{[EMNLP'26]}] K. Guo, X. Dai, \textbf{S. Zeng}, H. Shomer, H. Han, Y. Wang, J. Tang, ``Beyond Static Retrieval: Opportunities and Pitfalls of Iterative Retrieval in GraphRAG'', \textit{Conference on Empirical Methods in Natural Language Processing}, 2026.
\item[{[ICML'26]}] X. Dai, K. Yang, C. Luo, \textbf{S. Zeng}, K. Guo, J. Tang, ``When Do Hallucinations Arise? A Graph Perspective on the Evolution of Path Reuse and Path Compression'', \textit{International Conference on Machine Learning}, 2026.
\item[{[KDD'26]}] K. Guo, X. Dai, Z. Zhang, N. Lin, \textbf{S. Zeng}, J. Ren, H. Han, J. Tang, ``Why Retrieval-Augmented Generation Fails: A Graph Perspective'', \textit{ACM SIGKDD Conference on Knowledge Discovery and Data Mining}, 2026.
\item[{[NeurIPS'25]}] J. Ren, Z. Dai, X. Tang, Y. Xing, \textbf{S. Zeng}, H. Liu, J. Zeng, Q. Peng, S. Varshney, S. Wang, Q. He, C. C. Aggarwal, H. Liu, ``Keeping an Eye on LLM Unlearning: The Hidden Risk and Remedy'', \textit{Conference on Neural Information Processing Systems}, 2025.
\item[{[EMNLP'25]}] K. Guo, H. Shomer, \textbf{S. Zeng}, H. Han, Y. Wang, J. Tang, ``Empowering GraphRAG with Knowledge Filtering and Integration'', \textit{Conference on Empirical Methods in Natural Language Processing}, 2025.
\item[{[ACL'25]}] B. Wang, W. He, P. He, \textbf{S. Zeng}, Z. Xiang, Y. Xing, J. Tang, ``Unveiling Privacy Risks in LLM agent memory'', \textit{Annual Meeting of the Association for Computational Linguistics}, 2025.
\item[{[KDD'25]}] B. Li, Z. Chen, H. Han, \textbf{S. Zeng}, J. Liu, J. Tang, ``Unveiling Mode Connectivity in Graph Neural Networks'', \textit{ACM SIGKDD Conference on Knowledge Discovery and Data Mining}, 2025.
\item[{[CVPR'25]}] J. Ren, K. Chen, Y. Cui, \textbf{S. Zeng}, H. Liu, Y. Xing, J. Tang, L. Lyu, ``Six-CD: Benchmarking Concept Removals for Benign Text-to-Image Diffusion Models'', \textit{IEEE/CVF Conference on Computer Vision and Pattern Recognition}, 2025.
\item[{[TMLR]}] P. He, Y. Xing, H. Xu, J. Ren, Y. Cui, \textbf{S. Zeng}, J. Tang, M. Yamada, M. Sabokrou, ``Stealthy Backdoor Attack via Confidence-driven Sampling'', \textit{Transactions on Machine Learning Research}.
\item[{[ECCV'24]}] J. Ren, Y. Li, \textbf{S. Zeng}, H. Xu, L. Lyu, Y. Xing, J. Tang, ``Unveiling and Mitigating Memorization in Text-to-Image Diffusion Models through Cross Attention'', \textit{European Conference on Computer Vision}, 2024.
\item[{[EMNLP'24]}] H. Xu, J. Ren, P. He, Y. Cui, \textbf{S. Zeng}, H. Liu, J. Tang, A. Liu, ``On the Generalization of Training-based ChatGPT Detection Methods'', \textit{Findings of the Association for Computational Linguistics: EMNLP}, 2024.
\item[{[ICLR'24, Spotlight]}] P. He, H. Xu, J. Ren, Y. Cui, \textbf{S. Zeng}, H. Liu, C. Aggarwal, J. Tang, ``Sharpness-Aware Data Poisoning Attack'', \textit{International Conference on Learning Representations, \textbf{\color{honor}Spotlight}}, 2024.
\item[{[NeurIPS'23]}] J. Li, H. Shomer, H. Mao, \textbf{S. Zeng}, Y. Ma, N. Shah, J. Tang, D. Yin, ``Evaluating Graph Neural Networks for Link Prediction: Current Pitfalls and New Benchmarking'', \textit{Conference on Neural Information Processing Systems, Datasets \& Benchmarks Track}, 2023.
\end{publist}

\section{PREPRINTS}
    % \item \textbf{Shenglai Zeng}, Jiankun Zhang, Bingheng Li, Yuping Lin, Tianqi Zheng, Dante Everaert, Hanqing Lu, Hui Liu, Yue Xing, Monica Xiao Cheng, Jiliang Tang \\\textbf{\textit{Towards Knowledge Checking in Retrieval-augmented Generation: A Representation Perspective}}\\ \textbf{Submitted to ACL ARR}
%     \item \textbf{Shenglai Zeng*}, Jiankun Zhang*, Kai Guo, Xinnan Dai, Hui Liu, Jiliang Tang, Yi Chang\\\textbf{\textit{Fix Before Search: Benchmarking Agentic Query Visual Pre-processing in
% Multimodal Retrieval-augmented Generation}}\\ \textbf{Submitted to ICML}
\begin{publist}
\item[{[Preprint]}] X. Long, Z. Chen, \textbf{S. Zeng}, S. Wang, K. Guo, J. Tang, ``MemTrace: Probing What Final Accuracy Misses in Long-Term Memory''.
\item[{[Preprint]}] X. Dai, C.-H. Lo, K. Guo, \textbf{S. Zeng}, D. Luo, J. Tang, ``Uncovering Graph Reasoning in Decoder-only Transformers with Circuit Tracing''.
\item[{[Under review]}] Z. Chen, J. Gu, J. Yin, X. Long, \textbf{S. Zeng}, X. Liu, K. Guo, K. Zhou, J. Tang, ``Exploring Cross-Scenario Generality of Agentic Memory Systems: Diagnostics and a Strong Baseline'', \textit{ACL Rolling Review}.
\item[{[Under review]}] X. Dai, K. Guo, C.-H. Lo, \textbf{S. Zeng}, J. Ding, D. Luo, S. Mukherjee, J. Tang, ``GraphGhosts: Tracing Reasoning Structures Behind Large Language Models'', \textit{ACL Rolling Review}.
\item[{[Preprint]}] Y. Lin, P. He, H. Zhao, \textbf{S. Zeng}, Z. Xiang, H. Liu, J. Tang, ``Structural Alignment Faking: Hiding Malicious Capabilities in Mixture-of-Experts LLMs''.
\item[{[Preprint]}] J. Ren, H. Xu, P. He, Y. Cui, \textbf{S. Zeng}, J. Zhang, H. Wen, J. Ding, H. Liu, Y. Chang, J. Tang, ``Copyright Protection in Generative AI: A Technical Perspective''.
\end{publist}





% \section{KEY SKILLS}
% \begin{tabular}{ @{} >{\bfseries}l @{\hspace{6ex}} l }
% Programming Language\ & Python, C, C++, Java,Matlab \\
% Research Tool\ & Latex, Overleaf \\
% AI Framework\ & Pytorch, Mxnet,Tensorflow\\
% Network \ &  Cisco Certified Network Associate(CCNA)\\
% \end{tabular}

%\section{COMPETITIVE PROGRAMMING}
%{\customcvproject{Problem Solving}{}
%{\begin{itemize}
  %\item UVA online judge (80+ problems)
%\end{itemize}
%}
%}











% \section{SELECTED ACADEMIC PROJECTS}

% {\customcvproject{\textbf{Exploring Privacy Issues in Retrieval-Augmented Generation (RAG)}}{Oct 2023- Feb 2024}
%   {\begin{itemize}
%     \item We've uncovered two pivotal aspects: (1). Privacy challenges within RAG's own data (2). RAG's potential to safeguard training data 
%     \item \textbf{Data Leak Quantified:} RAG systems can leak private retrieval data, with our study showing about 50\% of sensitive retrieval data being output.
%     \item \textbf{Mitigation Efforts:} We've explored naive defenses such as summarization and retrieval thresholds. These methods help mitigate risks but don't completely resolve the issue, indicating the gravity of privacy risks in RAG.
%     \item  \textbf{Training Data Safeguard:} RAG shows promise in protecting training data, offering a strategy to bolster privacy in AI systems. 
%     \item Our code is available at \href{https://github.com/phycholosogy/RAG-privacy}{https://github.com/phycholosogy/RAG-privacy}
% \end{itemize}
%   }
  
% {\customcvproject{\textbf{Exploring Memorization in Fine-tuned Language Models}}{May 2023- Oct 2023}
%   {\begin{itemize}
%     \item Extensively studied the memorization effect of LLMs during the fine-tuning stage across different tasks.
%     \item \textbf{Fine-tuning Risks:} Utilizing copyrighted or private data in fine-tuning poses privacy/IP risks.
%     \item \textbf{Task Disparity:} Summarization \& Dialogue show high memorization, while QA and Classification are lower.
%     \item \textbf{Task-specific Scaling:} For high memorization tasks, memorization increases with larger models. Conversely, for low memorization tasks, increasing model size has little impact on memorization.
%     \item \textbf{Attention’s Role:} High memorization tasks have uniform, sparse attention patterns. We unravel the nuances between attention \& memorization. 
%     \item \textbf{Solution in Sight:} Multi-task fine-tuning buffers against high memorization vs single-task techniques.
% \end{itemize}
%   }


% {\customcvproject{Federated Learning Framework Design in  Industrial Metaverse}{Feb 2022 – Present}
% {\begin{itemize}
%   \item \textbf{Background}: This work mainly focuses on the non-i.i.d streaming data collected by distributed edge devices in Industrial Metaverse.(e.g. OCR applications in industrial parks.) 
%   \item \textbf{Challenge}: data heterogeneity/limited band width/ catastrophic forgetting towards streaming data
%   \item A dynamic FL training paradigm designed for rapidly changing streaming data while eliminating data heterogeneity.
%   \item A knowledge maintained online learning method for FL to prevent catastrophic forgetting
%   \item Task decomposition in federated learning to reduce communication pressure on edge networks
%   \item This work is accepted by IEEE Transactions on Cloud Computing.
% \end{itemize}
% }
% }

% {\customcvproject{Semi-Supervised Federated Learning for IoT}{Jun 2021 – Sept 2022}
% {\begin{itemize}
%   \item This work aims to develop an applicable federated learning mechanism for IoT devices in semi-supervised settings, considering personalization and communication efficiency simultaneously.  
%   \item We introduced neural network pruning techniques into semi-supervised federated learning and a novel
% structure-aware collaborative distillation approach which can aggregate models with different structures.
%   \item This work is accepted  by  SDM 2023.
% \end{itemize}
% }
% }



% \section{ONGOING PROJECTS}
% {\customcvproject{Federated Learning Framework Design in  Industrial IoT}{Feb 2022 – Present}
% {\begin{itemize}
%   \item \textbf{Background}: This work mainly focuses on the non-i.i.d streaming data collected by distributed edge devices in Industrial IoT.(e.g. OCR applications in industrial parks.) 
%   \item \textbf{Challenge}: data heterogeneity/limited band width/ catastrophic forgetting towards streaming data
%   \item A dynamic FL training paradigm designed for rapidly changing streaming data while eliminating data heterogeneity.
%   \item A knowledge maintained online learning method for FL to prevent catastrophic forgetting
%   \item Task decomposition in federated learning to reduce communication pressure on edge networks
%   \item This work is the extended version of FedGSP. As the first author, I will move forward with this work as soon as possible. 
% \end{itemize}
% }
% }
\section{Invited Talks}
\begin{minipage}{\maincolumnwidth}%
	\small{
    	\begin{itemize}
        			\item{``Securing Text and Empowering Vision: Evolution from RAG to Multimodal Agentic AI'' \hfill 2026 \\ Invited talk at University of Georgia, hosted by Prof.\ Zhen Xiang \emph{(upcoming)}}
        			\item{``Securing Text and Empowering Vision: Evolution from RAG to Multimodal Agentic AI'' \hfill 2026 \\ Invited talk at Washington University in St.\ Louis, hosted by Prof.\ Jiaxin Huang}
        			\item{``Attn-GS: Attention-Guided Context Compression for Efficient Personalized LLMs'' \hfill 2026 \\ Oral presentation at ACL 2026}
        			\item{``Towards Knowledge Checking in Retrieval-Augmented Generation: A Representation Perspective'' \hfill 2025 \\ Oral presentation at NAACL 2025}
        			\item{``Heterogeneous Federated Learning via Grouped Sequential-to-Parallel Training'' \hfill 2022 \\ Oral presentation at DASFAA 2022}
		\end{itemize}}%
\end{minipage}%

\section{Services}
\begin{minipage}{\maincolumnwidth}%
	\small{
    	\begin{itemize}
        			\item{\textbf{Conference Reviewer / Program Committee:} NeurIPS, ICML, ICLR; ACL, EMNLP, NAACL (via ACL Rolling Review), etc.}
           \item{\textbf{Journal Reviewer:} ACM Transactions on Information Systems (TOIS), IEEE TKDE, IEEE TMC, IEEE TPDS, IEEE Transactions on Information Forensics \& Security, IEEE Transactions on Services Computing, IEEE Transactions on Vehicular Technology, IEEE Internet of Things Journal, etc.}
           \item Reviewed \textbf{60+} papers.
        
        
          
        
		\end{itemize}}%
\end{minipage}%





% Publications from a BibTeX file without multibib
%  for numerical labels: \renewcommand{\bibliographyitemlabel}{\@biblabel{\arabic{enumiv}}}% CONSIDER MERGING WITH PREAMBLE PART
%  to redefine the heading string ("Publications"): \renewcommand{\refname}{Articles}
\section{Mentorship}
\begin{minipage}{\maincolumnwidth}%
	\small{
    	\begin{itemize}
        			\item{\textbf{Jiankun Zhang} (PhD@Jilin University): [ICML'26], [ACL'24]}
        			\item{\textbf{Xianxuan Long} (PhD@MSU): [NeurIPS'26], [Preprint]}
        			\item{\textbf{Bo Wang} (PhD@Jilin University): [ACL'25]}
        			\item{\textbf{Qirui Wang} (MS@Xi'an Jiaotong University): [NeurIPS'26]}
		\end{itemize}}%
\end{minipage}%

\nocite{*}
\bibliographystyle{plain}
\bibliography{publications}                        % 'publications' is the name of a BibTeX file

% Publications from a BibTeX file using the multibib package
%\section{Publications}
%\nocitebook{book1,book2}
%\bibliographystylebook{plain}
%\bibliographybook{publications}                   % 'publications' is the name of a BibTeX file
%\nocitemisc{misc1,misc2,misc3}
%\bibliographystylemisc{plain}
%\bibliographymisc{publications}                   % 'publications' is the name of a BibTeX file

%-----       letter       ---------------------------------------------------------

\end{document}


%% end of file `template.tex'.
